\documentclass[10pt,a4paper]{amsart}
\usepackage[margin=19mm]{geometry}
\usepackage{amsmath,amssymb,mathtools}
\usepackage[scr=rsfs,cal=euler]{mathalfa}
\usepackage{xfrac}
\usepackage[unicode,hidelinks,bookmarks=false]{hyperref}
\newcommand{\ie}{i.e.\ }
\newcommand{\cf}{cf.\ }
\newcommand{\resp}{respectively\ }
\newcommand{\Subsection}{\subsection}
\providecommand{\nobreakdash}{\nobreak\hbox{-}}
\setlength{\emergencystretch}{2em}
\multlinegap0pt
\usepackage{amssymb}
\usepackage{mathtools}
\usepackage[shortlabels]{enumitem}
\usepackage{tikz}
\usepackage{float}
\newtheorem{lemma}{Lemma}
\newcommand{\dist}{\operatorname{dist}}
\title{Two Arc-Disjoint Hamiltonian Paths in Finite Two-Generated Abelian Cayley Digraphs}
\author{SangHyun Park}
\date{Source: arXiv:2605.27241v1 (CC BY 4.0)}
\begin{document}
\maketitle

\noindent\emph{Source and licence.} Two Arc-Disjoint Hamiltonian Paths in Finite Two-Generated Abelian Cayley Digraphs. Version arXiv:2605.27241v1 (CC BY 4.0), distributed by arXiv under the Creative Commons Attribution 4.0 International licence, http://creativecommons.org/licenses/by/4.0/, as recorded in the arXiv licence metadata for this exact version and shown on the abstract page https://arxiv.org/abs/2605.27241v1. Full licence notice: \texttt{licenses/arxiv-2605.27241-CC-BY-4.0.txt}.\par\smallskip

\noindent\textit{Editorial scope.} Counted unit: the article's lemma lem:endpoint-blocked-growth (source0.tex lines 1519--1528). The article's complete original proof of it is delivered with it (source0.tex lines 1530--1621, one \texttt{proof} environment of 92 lines). No mathematics is written, repaired or re-derived: the statement and the proof are the article's own lines, in the article's own notation, and no retained line is substituted or re-flowed.  Retained uncounted as the source-local context the proof needs: lines 577--582; lines 744--746; lines 1273--1283; lines 1365--1375; lines 1404--1426; lines 1486--1495.  Nothing was omitted from any retained span, so the package carries no omission record.  No retained line contains a citation, so the wrapper carries no bibliography and no bibliography entry.  No retained line contains a figure environment, an image-inclusion command or drawing code, so no image is delivered and the package declares no asset dependency.  Editorial formatting only, in addition: an AMS \texttt{amsart} preamble that restates the article's own declarations and macros used by the retained lines, namely the environments \texttt{lemma} on separate counters, together with the standard packages those lines need.  The retained environments print in this document as Lemma 1 in order of appearance, whereas the article prints the same items with the numbering of its own full document.  The complete native source of the article as distributed in the arXiv e-print for this exact version is bundled unchanged as \texttt{sources/201478/source0.tex}.
\begin{lemma}\label{lem:theta-bound}
For positive integers $p,q$,
\[
        \Theta(p,q)\ge \frac{(p-1)(q-1)}2.
\]
\end{lemma}

\begin{lemma}[Location of reflected points]\label{lem:target-placement}
Assume $\delta=\dist(Z,N-Z)\ge2$.  Let $x\in Z$ and put $T=N-x$.  Then $T$ lies either in the $\delta$-core of an internal gap of $Z$, or in the $\delta$-core of one of the two endpoint caps.
\end{lemma}

\begin{lemma}[Endpoint cap-sector filling for \(c_L=2\)]\label{lem:width-three}
Assume \(c_L=2\).  Let \(R\) be an internal primitive ray generator to the right of the left endpoint ray \(C\), and put \(\alpha=H(R)\).  In this normalization the endpoint cap consists of the two multiples \(C\) and \(2C\); the sector families \(C+tR\) and \(2C+tR\) give the terms used in the strengthened boundary count.  Then the mass between \(C\) and \(R\) satisfies
\[
        M(C,R)\ge \alpha-2.
\]
If in addition \(N=4\alpha-2\), then
\[
        M(C,R)\ge \alpha-1.
\]
The reflected statement holds when \(c_R=2\).
\end{lemma}

\begin{lemma}[Boundary equalities with internal sign gaps]\label{lem:boundary-created-internal-edge}
Assume \(\delta=\dist(Z,N-Z)\ge2\).  Let a point lying on the boundary of an internal \(\delta\)-core create an edge of \(\Gamma_\delta\) whose sign-propagation gaps are internal: either
\[
        z_p+z_t=N-\delta
\]
with right gaps at both endpoints internal, or
\[
        z_p+z_{t+1}=N+\delta
\]
with left gaps at both endpoints internal.  Then no such boundary equality can occur.
\end{lemma}

\begin{lemma}[Prefix rigidity in a one-sided endpoint case]\label{lem:endpoint-blocked-initial}
Suppose
\[
        z_i+z_s=N-\delta,
        \qquad i<s,
\]
and suppose the propagation at $z_i$ is internally blocked by $\alpha=\lambda_i>\delta$.  Let $C=c_R=N-z_s$, so $z_i=C-\delta$, and assume that
\[
        T=N-z_{i+1}=z_s+\delta-2\alpha
\]
lies in the $\delta$-core of an internal gap $[z_t,z_{t+1}]$.  Then
\[
        \delta\le C<\alpha,
        \qquad
        c_L=0,
\]
and every gap before the $\alpha$-gap has multiplicity one:
\[
        \lambda_0=\lambda_1=\cdots=\lambda_{i-1}=1,
        \qquad
        z_r=2r\quad(0\le r\le i).
\]
\end{lemma}

\begin{lemma}[Position of the first reflected point]\label{lem:endpoint-blocked-position}
Under the hypotheses of Lemma~\ref{lem:endpoint-blocked-initial}, the point $T=N-z_{i+1}$ does not lie in the $\delta$-core of the $\alpha$-gap itself.  It lies strictly inside the $\delta$-core of a later internal gap $[z_t,z_{t+1}]$, where $t>i+1$ and, with $p=i+1$,
\[
        T=z_t+\delta+2\ell,
        \qquad
        1\le \ell\le \gamma-\delta-1,
        \qquad
        \gamma=\lambda_t\ge\delta+2.
\]
\end{lemma}

\begin{lemma}[Second reflection inequalities]\label{lem:endpoint-blocked-growth}
Under the hypotheses and notation of Lemma~\ref{lem:endpoint-blocked-position}, the following two inequalities hold:
\begin{align}
        \gamma&\ge \delta+2+\frac{\delta(\alpha-1)}2,
        \label{eq:endpoint-growth-one}\\
        \alpha&\ge \delta+1+\Theta(\gamma,C).
        \label{eq:endpoint-growth-two}
\end{align}
They are incompatible.
\end{lemma}

\begin{proof}
Put $p=i+1$.  Here $p$ indexes the cut value $z_p=z_{i+1}$; the internal blocking gap remains $G_i=[z_i,z_{i+1}]$.  We now reflect the left endpoint $z_t$ of the $\gamma$-gap, rather than the original cut value $z_p$.  From $T=N-z_p=z_t+\delta+2\ell$ we obtain
\[
        N-z_t=z_p+\delta+2\ell.
\]
We first show that $z_p<N-z_t<z_t$.  The left inequality is clear.  For the right one, the mass between the $\alpha$-gap $G_i$ and the $\gamma$-gap $G_t$ is
\[
        M(G_i,G_t)=\sum_{r=i+1}^{t-1}\lambda_r=\sum_{r=p}^{t-1}\lambda_r=\frac{z_t-z_p}{2}.
\]
Sector filling gives
\[
        M(G_i,G_t)
        \ge \Theta(\alpha,\gamma)
        \ge \frac{(\alpha-1)(\gamma-1)}2
        \ge \frac{\delta(\gamma-1)}2.
\]
Since $\ell\le\gamma-\delta-1$, this is greater than $\ell+\delta/2$.  Hence $z_t-z_p>\delta+2\ell$, proving $N-z_t<z_t$.

By Lemma~\ref{lem:target-placement}, $N-z_t$ lies in the $\delta$-core of an internal gap $[z_r,z_{r+1}]$ with $p\le r<t$.  Write $\eta=\lambda_r$ and
\[
        N-z_t=z_r+\delta+2u,
        \qquad
        0\le u\le \eta-\delta.
\]
The boundary cases $u=0$ and $u=\eta-\delta$ create, respectively, an internal negative or positive edge whose sign-propagation gaps are internal.  Lemma~\ref{lem:boundary-created-internal-edge} excludes these boundary equalities.  Thus
\[
        1\le u\le\eta-\delta-1,
        \qquad
        \eta\ge\delta+2.
\]
Moreover,
\[
        z_r=z_p+2(\ell-u).
\]
If $r=p$, then the adjacent gaps $\alpha$ and $\eta$ are both large, impossible.  Hence $r>p$, so $1\le u<\ell$.  The mass between the $\alpha$-gap $G_i$ and the $\eta$-gap $G_r$ is
\[
        M(G_i,G_r)=\sum_{q=i+1}^{r-1}\lambda_q=\ell-u,
\]
and sector filling gives
\[
        \ell-u\ge\Theta(\alpha,\eta)
        \ge\frac{(\alpha-1)(\eta-1)}2
        \ge\frac{\delta(\alpha-1)}2.
\]
Using $u\ge1$ and $\ell\le\gamma-\delta-1$ gives \eqref{eq:endpoint-growth-one}.

The second inequality comes from the suffix after the $\gamma$-gap.  Since
\[
        z_t+\delta+2\ell=T=z_s+\delta-2\alpha,
\]
we have the suffix upper bound
\[
        M(G_t,\mathcal B_R)=\sum_{r=t+1}^{s-1}\lambda_r=\alpha+\ell-\gamma\le \alpha-\delta-1.
\]
Because $C\ge\delta>0$, endpoint sector filling applies to the $\gamma$-gap and the endpoint cap $C$, so
\[
        M(G_t,\mathcal B_R)\ge\Theta(\gamma,C).
\]
Combining the two bounds gives
\[
        \Theta(\gamma,C)\le \alpha-\delta-1,
\]
which is \eqref{eq:endpoint-growth-two}.

If $(\delta,C)\ne(2,2)$, then $C\equiv\delta\pmod2$ implies either $\delta\ge3$ or $\delta=2$ and $C\ge4$.  Combining \eqref{eq:endpoint-growth-one}, \eqref{eq:endpoint-growth-two}, and Lemma~\ref{lem:theta-bound} gives
\[
        \alpha\ge
        \delta+1+
        \frac{(\delta+1)(C-1)}2
        +R(\alpha-1),
        \qquad
        R=\frac{\delta(C-1)}4.
\]
Here $R>1$.  Rearranging yields
\[
        (R-1)\alpha
        \le
        R-
        \left(\delta+1+\frac{(\delta+1)(C-1)}2\right),
\]
This right-hand side is
\[
        -\left(\delta+1+\frac{(\delta+2)(C-1)}4\right),
\]
which is negative.  Since $R-1>0$ and $\alpha>0$, this is impossible.

It remains to treat the boundary alignment isolated in Lemma~\ref{lem:width-three}, namely $\delta=C=2$.  Then \eqref{eq:endpoint-growth-one} gives $\gamma\ge\alpha+3$.  The reflected endpoint estimate for the multiplicity-two cap gives
\[
        \sum_{r=t+1}^{s-1}\lambda_r\ge \gamma-2,
\]
while the suffix bound above gives the same sum at most $\alpha-3$.  Hence $\gamma\le\alpha-1$, a contradiction.
\end{proof}

\end{document}
